\chapter{Conclusions}
\label{chap:conclusions}

This thesis investigated the use of neural stochastic flows for the reconstruction of AGN photometric light curves, a problem better suited to a probabilistic treatment. The damped random walk, whose transition probability is known in closed form and which admits an exact Gaussian process treatment, provided a setting in which the method could be validated not only against held-out data but against the analytic properties of the process it was trained to learn.

On simulated light curves at low noise, at both a fixed damping timescale and a distribution of timescales drawn from a physical prior, the method recovered the analytic transition and reconstructed the light curves to within a few per cent of the Bayes-optimal Gaussian process, and this performance was found to be insensitive to the weighting of the training objective. The method therefore learns the underlying dynamics and interpolates the gaps well when the observations are precise. When the sampling and noise were matched to those of the Zwicky Transient Facility, the point reconstruction remained close to the Gaussian process, but the predictive intervals became overconfident, and this miscalibration deepened as the number of noisy observations per curve increased. Applied to real ZTF light curves through a masking experiment, the method predicted withheld observations competitively with a Gaussian process fitted to the data, and both improved on linear interpolation. Across all of these experiments a consistent picture emerged: the reconstruction works, but its calibration depends on how measurement noise is treated, and the fitted Gaussian process remains a strong baseline on real data.

The limitation identified in these experiments points directly to the natural next step. The method treats the observed magnitudes as if they were the true signal, with no road through which the individual measurement uncertainties can enter, and it is this that causes the reconstruction to anchor too tightly to noisy observations. A model that instead consumes the per-epoch uncertainties, treating each observation as a noisy measurement of an unobserved latent state and propagating that uncertainty into its predictions, would address the limitation directly. Such a latent formulation, in which the flow describes the evolution of the underlying signal and a separate term relates it to the noisy observations, is the principal direction left for future work.

A second direction, motivated by the widest gaps found in the $i$ band, is the extension to multiple bands. The variability of an AGN is correlated across wavelength, as discussed in Chapter \ref{chap:variability}, so that a gap in one band coincides with observations in another. A model that reconstructs the bands jointly could use this coordinated variability to constrain a gap in one band with the data available in the others, recovering signal precisely where single-band reconstruction is most strongly limited.